\section{Actual normalized Cartier ratio metrics retain their original weights}

\begin{definition}[True original normalization point over a normal affine open]
\label{def:actual-normalization-relative-point-lift}
\lean{EnvelopeProperty.IntegralSchemeNormalization.relativePointLiftOnNormalAffine}
\lean{EnvelopeProperty.IntegralSchemeNormalization.map_relativePointLiftOnNormalAffine}
\leanok
\uses{lem:actual-normalization-affine-isomorphism}
An original relative scheme-valued point supported in an actual
original nonempty normal affine open has a true original lift
through the actual normalization. The lift retains the entire
original point, residue valuation, and prescribed base point.
\end{definition}
\begin{proof}
The actual normalization is an isomorphism over this entire
original affine open by the preceding original affine theorem.
Use the original valued-point homeomorphisms for the actual open
restrictions and this genuine whole original isomorphism, and then
map the resulting actual point through the actual source open
inclusion. The actual whole point-map composition identity says
that projection gives exactly the original entire relative point.
\end{proof}

\begin{theorem}[Original normal model ratio weight and equality]
\label{thm:actual-normalized-cartier-ratio-model-weight}
\lean{EnvelopeProperty.IntegralSchemeNormalization.normalizedCartierRatioWeight}
\lean{EnvelopeProperty.IntegralSchemeNormalization.normalizedCartierRatioWeight_eq}
\leanok
\uses{def:actual-normalization-relative-point-lift,cor:actual-finite-type-scheme-normalization-proper,thm:actual-integral-scheme-normalization-normal-stalks,cor:actual-cartier-ratio-model-weight,scheme:proper-model-ideal-pullback}
Let an original integral proper model admit an actual locally
finite-type structure over ANY field, and let $I,J$ be its
original Cartier ideals. At any original point supported in a
normal original affine open and with nonzero original $I$ norm,
form the actual original ratio line and canonical generic section
on its genuine proper normalization. Its true negative-log model
norm is exactly the difference of the original intrinsic $I,J$
model logarithms at the original entire point.
\end{theorem}
\begin{proof}
The actual original normalization is finite, hence proper over the
same original model base, and every actual original normalization
stalk is a normal domain. Its dominant integral projection retains
both actual original Cartier ideals, without flatness. Use the
actual point lift just constructed. The actual ideal-pullback norm
and logarithm identities, together with its whole point projection
identity, retain both original intrinsic ideal functions. In
particular the nonzero original $I$ norm proves the actual pulled
$I$ norm is nonzero. Construct the actual original ratio generic
section of these two true Cartier lines on this same actual normal
proper model. The proved true Cartier ratio metric identity gives
its negative logarithmic norm as the two actual pulled ideal
logarithms' difference. Their actual geometric pullback identities
identify this with the original ideal difference. Thus neither
properness of the normalization nor metric compatibility was
introduced as an additional input.
\end{proof}
